% =============================================================
%  CS102-F26 Digital Design and Computer Architecture I
%  Lecture 06 -- Logic Gates (Unit 6)
%  Follows Chapter 3 of A Student's Computer Architecture Book I.
%
%  Build the slides and the print handout:  npm run materials -- lecture-06
% =============================================================
\documentclass{wosm-slides}
\usepackage{cs102-lectures}[2026/10/19]

\title[Lecture 06 \textperiodcentered\ Logic Gates]{Logic Gates}
\subtitle{Lecture 06: NOT, AND, OR, NAND, NOR, XOR, XNOR, and universal gate sets}

% ---------- Helpers for this lecture --------------------------
% Wider input spacing than the course default, so input labels fit.
\tikzset{gates/.append style={
  and gate/.append style={logic gate input sep=0.36cm},
  or gate/.append style={logic gate input sep=0.36cm},
  nand gate/.append style={logic gate input sep=0.36cm},
  nor gate/.append style={logic gate input sep=0.36cm},
  xor gate/.append style={logic gate input sep=0.36cm},
  xnor gate/.append style={logic gate input sep=0.36cm}}}

% One gate with labelled inputs: \gate{<type>}{1 or 2}.
\newcommand*\gate[2]{%
  \begin{tikzpicture}[gates, baseline=(g.center)]
    \ifnum#2=1
      \node[#1 gate] (g) {};
      \draw[wire] (g.input) -- ++(-0.35,0) node[left] {$x$};
    \else
      \node[#1 gate, inputs=nn] (g) {};
      \draw[wire] (g.input 1) -- ++(-0.35,0) node[left] {$x$};
      \draw[wire] (g.input 2) -- ++(-0.35,0) node[left] {$y$};
    \fi
    \draw[wire] (g.output) -- ++(0.35,0);
  \end{tikzpicture}}

% A gate with its name, expression and truth table, in one column:
% \gatecard{<type>}{<NAME>}{<expression>}{<outputs for xy = 00,01,10,11>}
\newcommand*\gatecard[4]{%
  \centering
  \gate{#1}{2}\par\smallskip
  \textbf{#2}\par
  {\small $#3$}\par\smallskip
  \small
  \begin{tabular}{@{}cc|c@{}}
    $x$ & $y$ & \\
    \midrule
    \gatecardrows#4\relax
  \end{tabular}}
\def\gatecardrows#1,#2,#3,#4\relax{%
  0 & 0 & #1 \\ 0 & 1 & #2 \\ 1 & 0 & #3 \\ 1 & 1 & #4 \\}

% Proof steps: an equation and the law that justifies it.
\newenvironment{derivation}
  {\begin{tabular}{@{}r@{\;}l@{\qquad}l@{}}}
  {\end{tabular}}
\newcommand*\by[1]{\wosmlabel{#1}}

% =============================================================
\begin{document}

\maketitle

\begin{frame}{Today}
  \tableofcontents
\end{frame}

% =============================================================
\section{Logic gates}

\begin{frame}{What is a logic gate?}
  \begin{definition}[Logic gate]
    A \emph{logic gate} is a piece of hardware that computes one of the basic
    Boolean functions: addition, multiplication, negation, or a simple combination
    of these.
  \end{definition}

  \begin{itemize}
    \item Gates are connected by wires into \alert{circuits}, which compute larger
      functions.
    \item In a real chip, 0 and 1 are a low and a high voltage, and each gate is a
      handful of transistors: the switches of Lecture 01.
    \item From now on we work one layer up: with gates, not with transistors.
  \end{itemize}

  \medskip
  \wosmlabel{Every gate in this lecture is in the Gates library of
  Logisim-evolution.}
\end{frame}

\begin{frame}{NOT, AND and OR}
  \begin{columns}[T, totalwidth=\textwidth]
    \column{0.3\textwidth}
      \centering
      \gate{not}{1}\par\smallskip
      \textbf{NOT}\par
      {\small $x'$}\par\smallskip
      \small
      \begin{tabular}{@{}c|c@{}}
        $x$ & \\
        \midrule
        0 & 1 \\
        1 & 0 \\
      \end{tabular}
    \column{0.3\textwidth}
      \gatecard{and}{AND}{x \cdot y}{0,0,0,1}
    \column{0.3\textwidth}
      \gatecard{or}{OR}{x + y}{0,1,1,1}
  \end{columns}

  \medskip
  \begin{itemize}
    \item The three operations of Boolean algebra (Lecture 03), now as hardware.
    \item The small circle on the NOT gate, the \emph{bubble}, means
      \alert{complement}. Without it, the triangle is a \emph{buffer}: its output
      equals its input.
  \end{itemize}
\end{frame}

\begin{frame}{NAND and NOR}
  \begin{columns}[T, totalwidth=\textwidth]
    \column{0.3\textwidth}
      \gatecard{nand}{NAND}{(xy)' = x' + y'}{1,1,1,0}
    \column{0.3\textwidth}
      \gatecard{nor}{NOR}{(x + y)' = x'y'}{1,0,0,0}
    \column{0.34\textwidth}
      \begin{itemize}
        \item NAND is NOT-AND: an AND gate with a bubble on its output. NOR is
          NOT-OR.
        \item The second form of each expression is De Morgan's theorem
          (Lecture 03).
        \item Each one, on its own, can build every circuit: see the last section.
      \end{itemize}
  \end{columns}
\end{frame}

\begin{frame}{XOR and XNOR}
  \begin{columns}[T, totalwidth=\textwidth]
    \column{0.3\textwidth}
      \gatecard{xor}{XOR}{x \oplus y = x'y + xy'}{0,1,1,0}
    \column{0.3\textwidth}
      \gatecard{xnor}{XNOR}{(x \oplus y)' = x'y' + xy}{1,0,0,1}
    \column{0.34\textwidth}
      \begin{itemize}
        \item XOR, \emph{exclusive OR}: \enquote{either $x$ or $y$, but not both}.
          It is 1 when the inputs \alert{differ}.
        \item XNOR is NOT-XOR: 1 when the inputs are \alert{equal}.
        \item With more inputs, XOR is 1 when an odd number of them are 1.
      \end{itemize}
  \end{columns}
\end{frame}

\begin{frame}{All sixteen functions of two variables}
  \footnotesize
  \begin{center}
    \begin{tabular}{@{}cccc|l@{\qquad}cccc|l@{}}
      \multicolumn{4}{c|}{$xy = $ 00, 01, 10, 11} & &
      \multicolumn{4}{c|}{$xy = $ 00, 01, 10, 11} & \\
      \midrule
      0 & 0 & 0 & 0 & $0$ & 1 & 0 & 0 & 0 & $(x + y)'$ \quad NOR \\
      0 & 0 & 0 & 1 & $xy$ \quad AND & 1 & 0 & 0 & 1 & $x'y' + xy$ \quad XNOR \\
      0 & 0 & 1 & 0 & $xy'$ & 1 & 0 & 1 & 0 & $y'$ \\
      0 & 0 & 1 & 1 & $x$ & 1 & 0 & 1 & 1 & $x + y'$ \\
      0 & 1 & 0 & 0 & $x'y$ & 1 & 1 & 0 & 0 & $x'$ \\
      0 & 1 & 0 & 1 & $y$ & 1 & 1 & 0 & 1 & $x' + y$ \\
      0 & 1 & 1 & 0 & $x'y + xy'$ \quad XOR & 1 & 1 & 1 & 0 & $(xy)'$ \quad NAND \\
      0 & 1 & 1 & 1 & $x + y$ \quad OR & 1 & 1 & 1 & 1 & $1$ \\
    \end{tabular}
  \end{center}

  \normalsize
  \medskip
  These are the $2^{2^2} = 16$ functions counted in Lecture 04. Six of them are the
  two-input gates of this lecture, in complementary pairs: AND and NAND, OR and NOR,
  XOR and XNOR.
\end{frame}

\begin{frame}{Gates with more inputs}
  \begin{itemize}
    \item An $n$-input AND is 1 when \emph{all} its inputs are 1; an $n$-input OR
      when \emph{at least one} is.
    \item Because AND and OR are associative, a 3-input AND is the same as two
      2-input ANDs in a row: $xyz = (xy)z$.
    \item A 3-input NAND is $(xyz)'$, a NAND with three inputs.
  \end{itemize}

  \begin{warning}[NAND and NOR are not associative]
    Two 2-input NANDs in a row are not a 3-input NAND. With $x = y = 0$, $z = 1$:
    $\mathrm{NAND}(\mathrm{NAND}(x, y), z) = \mathrm{NAND}(1, 1) = 0$, but
    $(xyz)' = 1$.
  \end{warning}
\end{frame}

% =============================================================
\section{From expressions to circuits}

\begin{frame}{Drawing a circuit}
  Build $f = xy' + x'z$, the result of Lecture 05.

  \medskip
  \begin{columns}[c, totalwidth=\textwidth]
    \column{0.55\textwidth}
      \centering
      \begin{tikzpicture}[gates]
        \node (y) at (0,1.3) {$y$};
        \node (x) at (0,0.1) {$x$};
        \node (z) at (0,-1.4) {$z$};
        \node[not gate] (ny) at (1.4,1.3) {};
        \node[not gate] (nx) at (1.4,-0.6) {};
        \node[and gate, inputs=nn] (a1) at (3.0,0.95) {};
        \node[and gate, inputs=nn] (a2) at (3.0,-0.85) {};
        \node[or gate, inputs=nn] (o) at (4.6,0.05) {};
        \draw[wire] (y) -- (ny.input);
        \draw[wire] (ny.output) -- ++(0.25,0) |- (a1.input 1);
        \draw[wire] (x) -- (0.6,0.1) |- (a1.input 2);
        \draw[wire] (0.6,0.1) |- (nx.input);
        \fill[wosm-text] (0.6,0.1) circle (1.5pt);
        \draw[wire] (nx.output) -- ++(0.25,0) |- (a2.input 1);
        \draw[wire] (z) -- (2.35,-1.4) |- (a2.input 2);
        \draw[wire] (a1.output) -- ++(0.25,0) |- (o.input 1);
        \draw[wire] (a2.output) -- ++(0.25,0) |- (o.input 2);
        \draw[wire] (o.output) -- ++(0.4,0) node[right] {$f$};
        \node[edgelabel, above] at (a1.north) {$xy'$};
        \node[edgelabel, below] at (a2.south) {$x'z$};
      \end{tikzpicture}
    \column{0.41\textwidth}
      \begin{itemize}
        \item One gate per operation: two NOTs, two ANDs, one OR.
        \item A dot marks where one wire branches into two; crossing wires without a
          dot are not connected.
        \item A sum of products always gives this shape: NOTs, a column of ANDs,
          and one OR.
      \end{itemize}
  \end{columns}
\end{frame}

\begin{frame}{Building and testing in Logisim-evolution}
  \begin{enumerate}
    \item Place an input pin for each variable and an output pin for the result.
    \item Add the gates from the Gates library and connect them with wires.
    \item With the poke tool, set the inputs to every combination and check the
      output against the truth table.
  \end{enumerate}

  \bigskip
  \begin{remark}
    Testing every combination is proof by exhaustion again (Lecture 03). With 3
    inputs it is 8 checks; for larger circuits, Logisim-evolution can list the truth
    table for you (Project $\rightarrow$ Analyze Circuit).
  \end{remark}
\end{frame}

% =============================================================
\section{Universal gate sets}

\begin{frame}{Universality}
  \begin{definition}[Universal set of gates]
    A set of gates is \emph{universal} if every Boolean function can be built from
    its gates alone.
  \end{definition}

  \begin{theorem}
    $\{\mathrm{NOT}, \mathrm{AND}, \mathrm{OR}\}$ is universal.
  \end{theorem}

  \emph{Proof.} Every Boolean function has a canonical sum of products (Lecture 04),
  which uses only NOT, AND and OR.
\end{frame}

\begin{frame}{NOT and AND are enough}
  \begin{theorem}
    $\{\mathrm{NOT}, \mathrm{AND}\}$ is universal.
  \end{theorem}

  To show a set is universal, build the gates of a set already known to be
  universal. Here only OR is missing:

  \medskip
  \begin{columns}[c, totalwidth=\textwidth]
    \column{0.52\textwidth}
      \small
      \begin{tabular}{@{}r@{\;}l@{}}
        $x + y$ & $= ((x + y)')'$ \quad \by{involution} \\
                & $= (x'y')'$ \quad \by{De Morgan} \\
                & $= \mathrm{NOT}(\mathrm{AND}(\mathrm{NOT}\,x, \mathrm{NOT}\,y))$ \\
      \end{tabular}
    \column{0.44\textwidth}
      \centering
      \begin{tikzpicture}[gates]
        \node[not gate] (nx) at (0,0.4) {};
        \node[not gate] (ny) at (0,-0.4) {};
        \node[and gate, inputs=nn] (a) at (1.5,0) {};
        \node[not gate] (n) at (2.9,0) {};
        \draw[wire] (nx.input) -- ++(-0.3,0) node[left] {$x$};
        \draw[wire] (ny.input) -- ++(-0.3,0) node[left] {$y$};
        \draw[wire] (nx.output) -- ++(0.2,0) |- (a.input 1);
        \draw[wire] (ny.output) -- ++(0.2,0) |- (a.input 2);
        \draw[wire] (a.output) -- (n.input);
        \draw[wire] (n.output) -- ++(0.35,0) node[right] {$x + y$};
      \end{tikzpicture}
  \end{columns}

  \medskip
  By duality, $\{\mathrm{NOT}, \mathrm{OR}\}$ is universal too: $xy = (x' + y')'$.
\end{frame}

\begin{frame}{NAND alone is enough}
  \begin{theorem}
    $\{\mathrm{NAND}\}$ is universal.
  \end{theorem}

  \begin{columns}[T, totalwidth=\textwidth]
    \column{0.25\textwidth}
      \centering
      \begin{tikzpicture}[gates]
        \node[nand gate, inputs=nn] (g) at (0,0) {};
        \draw[wire] (g.input 1) -- ++(-0.3,0) coordinate (j) -- (j |- g.input 2)
          -- (g.input 2);
        \draw[wire] (j |- g.center) -- ++(-0.35,0) node[left] {$x$};
        \fill[wosm-text] (j |- g.center) circle (1.5pt);
        \draw[wire] (g.output) -- ++(0.3,0) node[right] {$x'$};
      \end{tikzpicture}\par\smallskip
      \small NOT: $\mathrm{NAND}(x, x) = x'$
    \column{0.36\textwidth}
      \centering
      \begin{tikzpicture}[gates]
        \node[nand gate, inputs=nn] (g1) at (0,0) {};
        \node[nand gate, inputs=nn] (g2) at (1.6,0) {};
        \draw[wire] (g1.input 1) -- ++(-0.3,0) node[left] {$x$};
        \draw[wire] (g1.input 2) -- ++(-0.3,0) node[left] {$y$};
        \draw[wire] (g1.output) -- ++(0.2,0) coordinate (j);
        \draw[wire] (j) |- (g2.input 1) (j) |- (g2.input 2);
        \fill[wosm-text] (j) circle (1.5pt);
        \draw[wire] (g2.output) -- ++(0.3,0) node[right] {$xy$};
      \end{tikzpicture}\par\smallskip
      \small AND: a NAND, then a NAND as NOT
    \column{0.37\textwidth}
      \centering
      \begin{tikzpicture}[gates]
        \node[nand gate, inputs=nn] (a) at (0,0.6) {};
        \node[nand gate, inputs=nn] (b) at (0,-0.6) {};
        \node[nand gate, inputs=nn] (g) at (1.35,0) {};
        \foreach \n/\v in {a/x, b/y} {
          \draw[wire] (\n.input 1) -- ++(-0.2,0) coordinate (j\n)
            -- (j\n |- \n.input 2) -- (\n.input 2);
          \draw[wire] (j\n |- \n.center) -- ++(-0.25,0) node[left] {$\v$};
        }
        \draw[wire] (a.output) -- ++(0.15,0) |- (g.input 1);
        \draw[wire] (b.output) -- ++(0.15,0) |- (g.input 2);
        \draw[wire] (g.output) -- ++(0.3,0) node[right] {$x + y$};
      \end{tikzpicture}\par\smallskip
      \small OR: $\mathrm{NAND}(x', y') = x + y$
  \end{columns}

  \bigskip
  NOT and AND are enough (previous slide), so NAND is enough.
\end{frame}

\begin{frame}{NOR alone is enough}
  \begin{theorem}
    $\{\mathrm{NOR}\}$ is universal.
  \end{theorem}

  \begin{proof}
    \small
    The duals of the NAND constructions. NOT: $\mathrm{NOR}(x, x) = (x + x)' = x'$.\par
    OR: $\mathrm{NOR}(\mathrm{NOR}(x, y), \mathrm{NOR}(x, y)) = ((x + y)')' = x + y$.\par
    AND: $\mathrm{NOR}(x', y') = (x' + y')' = xy$.\par
    NOT and OR are enough, so NOR is enough.
  \end{proof}

  \wosmlabel{The book builds a NAND from NOR gates instead: $xy = \mathrm{NOR}(x', y')$,
  then complement.}
\end{frame}

\begin{frame}{A set that is not universal}
  \begin{itemize}
    \item $\{\mathrm{AND}, \mathrm{OR}\}$ is \emph{not} universal: it cannot build
      NOT.
    \item Why: in a circuit of ANDs and ORs, changing an input from 0 to 1 can never
      change the output from 1 to 0. Each gate has this property, so the whole
      circuit has it.
    \item NOT does the opposite: as $x$ goes from 0 to 1, $x'$ goes from 1 to 0.
  \end{itemize}

  \bigskip
  \begin{remark}
    So the complement is what makes NAND and NOR special: each one combines an
    AND or an OR with a NOT.
  \end{remark}
\end{frame}

\begin{frame}{Why NAND and NOR matter}
  \begin{itemize}
    \item In CMOS, today's main chip technology, NOT takes 2 transistors, NAND and
      NOR take 4, and AND and OR take 6: a NAND or a NOR followed by a NOT.
    \item So NAND and NOR are the cheapest gates, and one gate type is enough for a
      whole chip.
  \end{itemize}

  \medskip
  \begin{columns}[c, totalwidth=\textwidth]
    \column{0.48\textwidth}
      \centering
      \begin{tikzpicture}[gates]
        \node[nand gate, inputs=nn] (a) at (0,0.65) {};
        \node[nand gate, inputs=nn] (b) at (0,-0.65) {};
        \node[nand gate, inputs=nn] (g) at (1.7,0) {};
        \draw[wire] (a.input 1) -- ++(-0.3,0) node[left] {$a$};
        \draw[wire] (a.input 2) -- ++(-0.3,0) node[left] {$b$};
        \draw[wire] (b.input 1) -- ++(-0.3,0) node[left] {$c$};
        \draw[wire] (b.input 2) -- ++(-0.3,0) node[left] {$d$};
        \draw[wire] (a.output) -- ++(0.2,0) |- (g.input 1);
        \draw[wire] (b.output) -- ++(0.2,0) |- (g.input 2);
        \draw[wire] (g.output) -- ++(0.3,0) node[right] {$ab + cd$};
      \end{tikzpicture}
    \column{0.48\textwidth}
      A sum of products becomes NAND--NAND: replace every gate by a NAND.
      \[ ((ab)' \cdot (cd)')' = ab + cd \]
      \wosmlabel{De Morgan, then involution.}
  \end{columns}
\end{frame}

\begin{frame}{Try it}
  \begin{enumerate}
    \item Build XOR from four NAND gates.\par
      \uncover<2->{\small $n = \mathrm{NAND}(x, y)$, \quad
      $x \oplus y = \mathrm{NAND}(\mathrm{NAND}(x, n), \mathrm{NAND}(y, n))$.
      Check it on all four inputs.}
    \item Build NOR from NAND gates only. How many do you need?\par
      \uncover<3->{\small Four: OR takes three (two as NOTs, one NAND), then one
      more as NOT.}
    \item Is $\{\mathrm{XOR}\}$ universal?\par
      \uncover<4->{\small No. A circuit of XORs only ever computes the XOR of some
      of its inputs (or 0), and $xy$ is not of that form.}
  \end{enumerate}
\end{frame}

% =============================================================
\section{Summary}

\begin{frame}{Summary}
  \begin{itemize}
    \item A logic gate computes a basic Boolean function in hardware: NOT, AND, OR,
      and the combinations NAND, NOR, XOR, XNOR.
    \item A bubble means complement. XOR is 1 when its inputs differ; XNOR when they
      are equal.
    \item A circuit has one gate per operation; a sum of products gives NOTs, ANDs,
      then an OR.
    \item A set of gates is universal if it can build every Boolean function:
      $\{\mathrm{NOT}, \mathrm{AND}, \mathrm{OR}\}$, $\{\mathrm{NOT}, \mathrm{AND}\}$,
      $\{\mathrm{NAND}\}$ and $\{\mathrm{NOR}\}$ are; $\{\mathrm{AND}, \mathrm{OR}\}$
      is not.
  \end{itemize}
\end{frame}

\begin{frame}{Before the next lecture}
  \begin{itemize}
    \item Read Chapter~3 of the book, \emph{Logic Gates \& Circuits}.
    \item In Logisim-evolution, build NOT, AND and OR from NAND gates only, and test
      each one on every input.
    \item Build $f = xy' + x'z$ with NOT, AND and OR gates, and check it against its
      truth table.
  \end{itemize}

  \begin{block}{Next lecture \unitref{7}}
    Combinational circuits: multiplexers and demultiplexers.
  \end{block}
\end{frame}

\end{document}
